\documentclass[11pt]{article}
\usepackage[a4paper,margin=1in]{geometry}
\usepackage{times}
\usepackage{amsmath,amssymb}
\usepackage{hyperref}
\usepackage{enumitem}
\hypersetup{hidelinks}

\begin{document}

\begin{center}
\Large\textbf{Statement on Prior Publication}
\end{center}

\vspace{1em}

\noindent\textbf{Submitted manuscript:} Beyond English-Only GRPO:
Training Language and Auxiliary Reward as Implicit Regularizers in
Sub-3B Mathematical Reasoning

\noindent\textbf{Author:} Vu Dang (Independent Researcher, Vietnam;
ORCID: \href{https://orcid.org/0009-0005-2344-1030}{0009-0005-2344-1030})

\vspace{1em}\hrule\vspace{1em}

\subsection*{Prior preprint disclosed}

A preliminary single-seed snapshot of this work was self-archived on
Zenodo on 7~May~2026 (DOI:
\href{https://doi.org/10.5281/zenodo.20061328}%
{10.5281/zenodo.20061328}). The preprint has not undergone peer review.
The Zenodo deposit will be updated to a new version pointing to the
present manuscript so that the deposit and submission remain in sync.

The accompanying open-source software repository
(\href{https://github.com/vudang4494/xling-grpo-sub3b}%
{github.com/vudang4494/xling-grpo-sub3b}, Apache-2.0) is shared between
the original Zenodo deposit and the present submission; the repository
contains code, configurations, model checkpoints, and evaluation
JSONs for the manuscript as submitted.

\vspace{1em}

\subsection*{Extensions over the original Zenodo deposit}

The submitted manuscript extends the original deposit with the
following new experimental content, methodological additions, and
revised conclusions:

\begin{enumerate}[itemsep=4pt]

\item \textbf{Multi-seed verification.}
The original deposit reported single-seed results (\texttt{seed=42})
for three training arms. The submitted manuscript expands to three
random seeds (\texttt{seed} $\in \{42, 123, 7\}$) per arm. This
reveals that the original single-seed claims (e.g., $+7.5$\,pp lift
on AMC-23, $+13.3$\,pp recovery on AIME-2024 via $R_5$) fall within
seed variance ($\sigma = 11.3$\,pp on AMC-23) and have been replaced
with multi-seed mean$\,\pm\,$standard deviation and bootstrap
$95\,\%$ confidence intervals for every quantitative claim.

\item \textbf{Mechanism ablation (new arm).}
The submitted manuscript adds Arm A4 (constant-bias reward
$R_{\text{const}}=1.0$, three seeds) to test the reward-magnitude
mechanism hypothesis. The ablation refutes a pure reward-magnitude
explanation of the $R_5$ effect on AIME-2024 \texttt{maj@8}:
$\armC{}-\armD{} = +5.6$\,pp, $95\,\%$ subject-bootstrap CI
$[+1.1, +11.1]$, excludes $0$.

\item \textbf{Multilingual evaluation (new benchmark).}
The submitted manuscript adds a 10-language MGSM evaluation across all
training arms plus the untrained base, comprising 80 new evaluation
JSONs (8 conditions $\times$ 10 languages). The result is a negative
finding: all training conditions converge to within $\pm 0.5$\,pp of
the untrained base mean.

\item \textbf{Reproducibility infrastructure.}
The submitted manuscript adds (a)~a SHA256-checksummed manifest
covering all training and evaluation artifacts, (b)~bootstrap
$95\,\%$ confidence intervals for all numerical claims, and (c)~a
per-component verification document (\texttt{VERIFICATION.md}).

\item \textbf{Reframed contributions.}
The submitted manuscript reframes around three orthogonal findings:
(i) positive --- A3 best on AIME-2024 \texttt{maj@8};
(ii) methodological --- multi-seed instability of vanilla EN GRPO;
(iii) negative --- no cross-lingual transfer on MGSM. The mechanism
hypothesis is presented as an open question rather than a confirmed
claim.

\end{enumerate}

\vspace{1em}

\subsection*{No other prior publication}

This manuscript has not been submitted to any peer-reviewed venue,
conference, or journal other than IEEE Access. No portions of it have
been presented at any conference or workshop. There are no other
related publications or preprints by the author covering this work.

\vspace{2em}\hrule\vspace{1em}

\noindent Signed:

\medskip

\noindent Vu Dang \\
\href{mailto:vu.dh4494@gmail.com}{vu.dh4494@gmail.com} \\
ORCID: \href{https://orcid.org/0009-0005-2344-1030}{0009-0005-2344-1030}

\end{document}
